\documentclass[11pt]{article}

\usepackage[a4paper, margin=1in]{geometry}
\usepackage{amsmath, amssymb, amsthm}
\usepackage{mathtools}
\usepackage{booktabs}
\usepackage{array}
\usepackage[round]{natbib}
\usepackage[colorlinks=true, linkcolor=blue, citecolor=blue, urlcolor=blue]{hyperref}
\usepackage{cleveref}
\usepackage{enumitem}
\usepackage{bm}

\newtheorem{theorem}{Theorem}[section]
\newtheorem{proposition}[theorem]{Proposition}
\newtheorem{lemma}[theorem]{Lemma}
\newtheorem{corollary}[theorem]{Corollary}
\newtheorem{definition}[theorem]{Definition}
\newtheorem{assumption}[theorem]{Assumption}
\newtheorem{example}[theorem]{Example}
\theoremstyle{remark}
\newtheorem{remark}[theorem]{Remark}

\DeclareMathOperator*{\argmax}{arg\,max}

\title{Quantitative Well-Posedness of Mean-Field Games\\ with Common Noise}
\author{Nisong Monyimba\thanks{Missouri University of Science and Technology. Email: \texttt{nisongmonyimba278@gmail.com}. Code and supplementary material: \url{https://github.com/NisongMonyimba/AppliedMathThesis}.}}
\date{\today}

\begin{document}
\maketitle

\begin{abstract}
We study mean-field games (MFGs) in which a continuum of agents interacts through a
common, shared Brownian motion in addition to idiosyncratic noise. The equilibrium is
characterised by a system of forward-backward stochastic differential equations
(FBSDEs) coupled through the conditional law of the representative agent's state given
the common-noise filtration. Building on the well-posedness theory of
\citet{carmonadelarue2018} and \citet{ahujarenyang2019}, we derive an explicit,
parameter-dependent contraction estimate for the associated fixed-point map in a
weighted Wasserstein-2 metric. The contraction constant depends explicitly on the
Lipschitz constant of the drift, the Frobenius norm of the common-noise coefficient,
and bounds on the relevant Lions derivatives, and is computable from the model
primitives. Existence and uniqueness of the mean-field equilibrium follow via the
Banach fixed-point theorem, and we verify the associated $\varepsilon$-Nash property
for the finite-player game. To the best of our knowledge, an explicit, computable
contraction bound of this form -- suitable as a design criterion for numerical
algorithms -- has not previously been derived in the common-noise setting; existing
work establishes existence and uniqueness but not quantitative stability. We illustrate
the theory on the canonical linear-quadratic MFG, for which all constants are explicit.
\end{abstract}

\section{Introduction}
\label{sec:intro}

Mean-field games (MFGs), introduced independently by \citet{lasrylions2007} and
\citet{huangetal2006}, model the equilibrium behaviour of a large population of
interacting agents, each optimising an individual objective that depends on the
aggregate behaviour of the population. In the classical formulation, agents are subject
only to idiosyncratic noise, and the mean-field limit is deterministic: the empirical
measure converges to a deterministic flow of measures $(\mu_t)_{t\in[0,T]}$.

Many applications, however, involve a source of randomness shared by \emph{all}
agents simultaneously -- a common macroeconomic shock, a shared environmental factor,
or systemic market risk. This motivates the study of \emph{common-noise} MFGs, in
which each agent's state is driven by both an idiosyncratic Brownian motion and a
shared Brownian motion $W^0$. In this setting the equilibrium measure flow
$\mu_t = \mathcal{L}(X_t \mid \mathcal{F}_t^{W^0})$ is itself a \emph{stochastic}
process, adapted to the common-noise filtration -- a genuinely more delicate object
than in the classical case.

The probabilistic well-posedness theory for common-noise MFGs was developed by
\citet{carmona2013probabilistic} and treated comprehensively by
\citet{carmonadelarue2018}, and extended under weaker (displacement) monotonicity
conditions by \citet{ahujarenyang2019}. These works establish existence and uniqueness
of the equilibrium under monotonicity-type structural conditions, but do not provide
\emph{explicit, parameter-dependent} contraction constants: quantities that tell a
practitioner, in terms of the model's Lipschitz constants and noise intensities,
precisely how fast (and under what conditions) a fixed-point iteration converges.

\paragraph{Contribution.} We derive such an explicit contraction estimate. Working
under standard Lipschitz, monotonicity, and regularity assumptions (\Cref{sec:setup}),
we show that the natural fixed-point map $\Phi$ associated with the common-noise MFG
equilibrium satisfies a Wasserstein stability estimate (\Cref{thm:stability}) whose
constants are given explicitly in terms of the drift Lipschitz constant $L_b$, the
common-noise coefficient norm $\|\sigma_0\|_F$, and bounds on the relevant Lions
derivatives. This yields a contraction in a weighted Wasserstein metric
(\Cref{thm:wellposedness}) with an explicit rate $\rho < 1$, from which existence,
uniqueness, and the $\varepsilon$-Nash property for the associated $N$-player game
follow. We illustrate all constants explicitly for the canonical linear-quadratic MFG.

We emphasise precisely what this contributes relative to the closely related recent
work of \citet{jacksontangpi2024} and \citet{jacksonmeszaros2025}, discussed in detail
below: those papers establish quantitative \emph{convergence rates of $N$-player Nash
equilibria to an assumed-to-exist mean-field equilibrium} via displacement
monotonicity, whereas the present paper constructs the mean-field equilibrium itself
as the fixed point of $\Phi$, via an explicit contraction estimate intended as a
design criterion for a numerical fixed-point iteration scheme. To the best of our
knowledge, relative to the specific works cited in this section, an explicit
computable Banach-contraction estimate of this form for the equilibrium
fixed-point map itself has not previously been derived in the common-noise setting.
This scope reflects the mean-field-game literature most directly
relevant to the present setting, including the papers identified in the focused
literature search underlying this claim, and is not a claim of an exhaustive survey
of all adjacent fields.

\paragraph{Related work.} \citet{cardaliaguetetal2019} derive the master equation for
common-noise MFGs and
establish PDE-level regularity; \citet{cardaliaguetseegersouganidis2025} resolve
further SPDE regularity questions in this setting. Neither addresses quantitative
FBSDE contraction rates for the fixed-point map $\Phi$. For classical
(idiosyncratic-only) MFGs,
\citet{cardaliaguet2010notes} gives a PDE treatment of well-posedness; numerical
schemes for this setting are studied by \citet{achdoucapuzzo2010} and
\citet{carmonalauriere2022}. The propagation-of-chaos theory underlying the
$\varepsilon$-Nash verification traces to \citet{mckean1966} and
\citet{sznitman1991}, with quantitative Wasserstein rates due to
\citet{fournierguillin2015}.

\textbf{Closely related work using synchronous coupling for common-noise FBSDE
stability.} \citet{jacksontangpi2024} use the same synchronous-coupling strategy
underlying \Cref{sec:stability} to derive quantitative \emph{convergence rates}
$O(1/N)$ of $N$-player open-loop Nash equilibria to a mean-field equilibrium, for
common-noise MFGs with controlled volatility, under a displacement-monotonicity
structural condition. This replaces the small-time-horizon or dissipativity
restriction used in earlier coupling-based arguments (e.g.\ the idiosyncratic-noise
result of Lauri\`ere and Tangpi that they build on) with a monotonicity condition
that holds more broadly. \citet{jacksonmeszaros2025} extend this to mean-field
games of control. Both works also establish existence and uniqueness of the
mean-field equilibrium, but via the \emph{method of continuation} for monotone
FBSDEs (in the sense of \citealp{peng1999}), not via a Banach-fixed-point
contraction estimate on the space of measure flows. The present paper's target
object -- an explicit contraction constant for the map $\Phi:\mu \mapsto \Phi(\mu)$
itself, usable directly as a step-size / convergence-rate criterion for a numerical
fixed-point iteration -- is accordingly a different (though methodologically
related) result. We view the displacement-monotonicity condition used by
\citet{jacksontangpi2024,jacksonmeszaros2025} as a promising route to relaxing the
smallness condition \eqref{eq:smallness-condition} required to close the Gronwall
argument in \Cref{app:algebra}, and we flag this as a direction for future work.
Finally, \citet{frizetal2026} study a linear-quadratic MFG with \emph{rough} (rather
than It\^o/Brownian) common noise and derive stability estimates with respect to
initial data and the common-noise path in that pathwise setting; the technical
setting (rough-path theory) and the notion of stability differ from the present
Brownian, expectation-level analysis.

\paragraph{Organisation.} \Cref{sec:setup} states the model, assumptions, and the
fixed-point formulation. \Cref{sec:stability} contains the core stability estimate.
\Cref{sec:main} states and proves the main theorem (contraction, existence,
uniqueness, $\varepsilon$-Nash). \Cref{sec:discussion} discusses the contraction
constant, its dependence on model parameters, and limitations of the current
approach. \Cref{sec:example} illustrates the theory on the linear-quadratic benchmark.

\section{Model and Standing Assumptions}
\label{sec:setup}

We work on a filtered probability space
$(\Omega, \mathcal{F}, (\mathcal{F}_t)_{t\in[0,T]}, \mathbb{P})$ supporting an
idiosyncratic Brownian motion $W$ and an independent common Brownian motion $W^0$.
Let $\mathcal{P}_2(\mathbb{R}^d)$ denote the space of Borel probability measures on
$\mathbb{R}^d$ with finite second moment, equipped with the Wasserstein-2 distance
$W_2$. The representative agent's state solves the McKean--Vlasov SDE
\begin{equation}
    dX_t^\mu = b(t, X_t^\mu, \mu_t, \alpha_t)\, dt + \sigma\, dW_t + \sigma_0\, dW_t^0,
    \qquad X_0^\mu \sim \mu_{\mathrm{init}},
    \label{eq:agent-sde}
\end{equation}
for a candidate measure flow $\mu = (\mu_t)_{t\in[0,T]}$, and the agent maximises
\begin{equation}
    J^\mu(\alpha) = \mathbb{E}\left[\int_0^T f(t, X_t^\mu, \mu_t, \alpha_t)\, dt +
    g(X_T^\mu, \mu_T)\right].
    \label{eq:objective}
\end{equation}

\begin{assumption}[Lipschitz drift]
\label{ass:lipschitz-drift}
There exists $L_b > 0$ such that for all $t, x, x', \mu, \mu', \alpha, \alpha'$,
\[
    \|b(t,x,\mu,\alpha) - b(t,x',\mu',\alpha')\| \le L_b\big(\|x-x'\| + W_2(\mu,\mu') +
    \|\alpha-\alpha'\|\big).
\]
\end{assumption}

\begin{assumption}[Lipschitz rewards]
\label{ass:lipschitz-rewards}
$f$ and $g$ are Lipschitz continuous in all arguments, with constants $L_f, L_g > 0$.
\end{assumption}

\begin{assumption}[Regularity and Lions differentiability]
\label{ass:regularity}
$b, f, g$ are $C^1$ in $x$ and Lions-differentiable in $\mu$ (in the sense of
\citealp[Ch.~5]{carmonadelarue2018}), with bounded, Lipschitz Lions derivatives.
\end{assumption}

\begin{assumption}[Bounded volatility]
\label{ass:volatility}
$\sigma > 0$ is constant, and $\sigma_0 \in \mathbb{R}^{d\times d_0}$ satisfies
$\|\sigma_0\|_F \le M_\sigma < \infty$.
\end{assumption}

\begin{assumption}[Uniform Hamiltonian concavity]
\label{ass:concavity}
$\alpha \mapsto H(t,x,\mu,p,\alpha) := p\cdot b(t,x,\mu,\alpha) + f(t,x,\mu,\alpha)$
is uniformly $\lambda_H$-concave for some $\lambda_H > 0$.
\end{assumption}

\begin{assumption}[Initial condition]
\label{ass:initial}
$\mu_{\mathrm{init}} \in \mathcal{P}_2(\mathbb{R}^d)$ has finite fourth moment.
\end{assumption}

Under \Cref{ass:lipschitz-drift,ass:lipschitz-rewards,ass:regularity,ass:concavity},
the Stochastic Maximum Principle characterises the optimal control for a \emph{fixed}
$\mu$: there is a unique solution $(X^\mu, p^\mu, q^\mu, r^\mu, \alpha^\mu)$ to
\begin{align}
    dX_t^\mu &= b(t, X_t^\mu, \mu_t, \alpha_t^\mu)\, dt + \sigma\, dW_t + \sigma_0\,
        dW_t^0, & X_0^\mu &\sim \mu_{\mathrm{init}}, \notag \\
    -dp_t^\mu &= \partial_x H(t, X_t^\mu, \mu_t, p_t^\mu, \alpha_t^\mu)\, dt -
        q_t^\mu\, dW_t - r_t^\mu\, dW_t^0, & p_T^\mu &= \partial_x g(X_T^\mu, \mu_T),
        \label{eq:agent-fbsde} \\
    \alpha_t^\mu &= \argmax_{\alpha} H(t, X_t^\mu, \mu_t, p_t^\mu, \alpha), \notag
\end{align}
following standard results for Lipschitz McKean--Vlasov FBSDEs
\citep[Vol.~I, Ch.~4]{carmonadelarue2018}. This defines the fixed-point map
\begin{equation}
    \Phi(\mu)_t := \mathcal{L}\big(X_t^\mu \mid \mathcal{F}_t^{W^0}\big), \qquad
    t \in [0,T],
    \label{eq:fixed-point-map}
\end{equation}
on the space of adapted measure flows $\mathcal{M} := D([0,T]; \mathcal{P}_2(\mathbb{R}^d))$
(the Skorokhod space of c\`adl\`ag measure-valued processes adapted to
$(\mathcal{F}_t^{W^0})$). A mean-field equilibrium is a fixed point $\Phi(\mu^*) = \mu^*$.

For $\lambda > 0$, define the weighted norm on $\mathcal{M}$,
\[
    \|\mu\|_{\lambda,T} := \sup_{t\in[0,T]} e^{-\lambda t}\,
    \mathbb{E}\big[W_2(\mu_t, \delta_0)^2\big]^{1/2},
\]
and the corresponding distance $\|\mu^1-\mu^2\|_{\lambda,T}$. This weighted norm
is metrically equivalent to the standard norm of \citet{carmonadelarue2018} but
absorbs the Gronwall growth rate more conveniently for extracting an explicit
contraction factor.

\section{Stability Estimate for the Fixed-Point Map}
\label{sec:stability}

The heart of the argument is a quantitative Lipschitz-type estimate for $\Phi$.

\subsection{Synchronous coupling}

Given two candidate measure flows $\mu^1, \mu^2 \in \mathcal{M}$, we couple the two
systems \eqref{eq:agent-fbsde} on a common probability space, using the \emph{same}
realisations of the idiosyncratic noise $W$ and, crucially, the same common noise
$W^0$. Since the Wasserstein distance between two laws is bounded above by the
$L^2$-distance between any coupling,
\begin{equation}
    \mathbb{E}\big[W_2(\Phi(\mu^1)_t, \Phi(\mu^2)_t)^2\big] \le
    \mathbb{E}\big[|X_t^1 - X_t^2|^2\big].
    \label{eq:wasserstein-l2-bound}
\end{equation}

\subsection{The difference process}

Write $\Delta X_t := X_t^1 - X_t^2$. Under synchronous coupling with identical,
constant diffusion coefficients, the stochastic terms cancel \emph{pathwise}:
\[
    d(\Delta X_t) = \big[b(t,X_t^1,\mu_t^1,\alpha_t^1) - b(t,X_t^2,\mu_t^2,\alpha_t^2)
    \big]\, dt,
\]
so that It\^o's formula applied to $|\Delta X_t|^2$ has no quadratic-variation
correction:
\begin{equation}
    d|\Delta X_t|^2 = 2\Delta X_t \cdot \big[b(t,X_t^1,\mu_t^1,\alpha_t^1) -
    b(t,X_t^2,\mu_t^2,\alpha_t^2)\big]\, dt.
    \label{eq:ito-expansion}
\end{equation}

Using \Cref{ass:lipschitz-drift} and Young's inequality,
\begin{equation}
    d|\Delta X_t|^2 \le \big[(2L_b + L_b^2)|\Delta X_t|^2 + L_b W_2(\mu_t^1,\mu_t^2)^2 +
    L_b|\alpha_t^1-\alpha_t^2|^2\big]\, dt.
    \label{eq:gronwall-setup}
\end{equation}

\subsection{Control difference via the SMP}

By the implicit function theorem applied to the optimality condition
$\partial_\alpha H = 0$, uniform concavity (\Cref{ass:concavity}) gives a constant
$C_\alpha > 0$ such that
\begin{equation}
    |\alpha_t^1 - \alpha_t^2| \le C_\alpha\big(|X_t^1-X_t^2| + |p_t^1-p_t^2| +
    W_2(\mu_t^1,\mu_t^2)\big).
    \label{eq:control-lipschitz}
\end{equation}
The Lipschitz dependence on the measure argument follows from the Lions
differentiability in \Cref{ass:regularity} together with
\citet[Prop.~5.35]{carmonadelarue2018}.

\subsection{Adjoint process estimate}

Differencing the adjoint BSDEs gives a linear BSDE for $\Delta p_t := p_t^1 - p_t^2$
with a mean-field coupling term controlled, via \Cref{ass:regularity}, by a constant
$L_\mu$ bounding the relevant Lions-derivative operator norm. Standard BSDE estimates
\citep{elkarouietal1997,zhang2017} yield
\begin{align}
    \mathbb{E}\bigg[\sup_{t\in[0,T]} |\Delta p_t|^2 + \int_0^T \big(|\Delta q_t|^2 +
    |\Delta r_t|^2\big)\, dt\bigg] \notag \\
    \le C_{\mathrm{BSDE}}\, \mathbb{E}\bigg[|\Delta X_T|^2 + \int_0^T \big(|\Delta
    X_t|^2 + |\alpha_t^1-\alpha_t^2|^2 + W_2(\mu_t^1,\mu_t^2)^2\big)\, dt\bigg],
    \label{eq:bsde-estimate}
\end{align}
where $C_{\mathrm{BSDE}}$ depends on the Lipschitz bounds, $T$, and -- crucially,
through the martingale representation terms driven by $W^0$ -- on $\|\sigma_0\|_F$.

\subsection{Closing the estimate}

Substituting \eqref{eq:control-lipschitz} into \eqref{eq:gronwall-setup} and
\eqref{eq:bsde-estimate} and setting
$y_t := \mathbb{E}[|\Delta X_t|^2] + \mathbb{E}[|\Delta p_t|^2]$, algebraic
manipulation (given in full in \Cref{app:algebra}) yields the closed differential
inequality
\begin{equation}
    \frac{d}{dt} y_t \le C_1\, y_t + C_2\, \mathbb{E}\big[W_2(\mu_t^1,\mu_t^2)^2\big],
    \label{eq:gronwall-differential}
\end{equation}
with
\begin{equation}
    C_1 = 2L_b^2 + 2\|\sigma_0\|_F^2 + C_{\mathrm{reg}}, \qquad
    C_2 = L_b + C_\alpha L_b + C_{\mathrm{BSDE}} L_\mu,
    \label{eq:constants}
\end{equation}
where $C_{\mathrm{reg}}$ collects the regularity constants from the Lions
derivatives and BSDE bounds. By Gronwall's inequality and
\eqref{eq:wasserstein-l2-bound}, this gives the main stability estimate:

\begin{theorem}[Wasserstein Stability]
\label{thm:stability}
Under Assumptions~\ref{ass:lipschitz-drift}--\ref{ass:initial}, for all
$\mu^1, \mu^2 \in \mathcal{M}$ and $t \in [0,T]$,
\begin{equation}
    \mathbb{E}\big[W_2(\Phi(\mu^1)_t, \Phi(\mu^2)_t)^2\big] \le C_2 \int_0^t
    e^{C_1(t-s)}\, \mathbb{E}\big[W_2(\mu_s^1,\mu_s^2)^2\big]\, ds,
    \label{eq:stability-estimate}
\end{equation}
with $C_1, C_2$ as in \eqref{eq:constants}.
\end{theorem}

\begin{remark}[Expectation, not pathwise]
The estimate \eqref{eq:stability-estimate} holds for the \emph{expected} squared
Wasserstein distance. The first-order stochastic terms driven by $W^0$ cancel in
expectation (via the martingale property) but not pathwise; this is why the
contraction below is stated in a weighted expectation norm rather than almost surely.
\end{remark}

\section{Main Theorem: Contraction, Existence, Uniqueness, and $\varepsilon$-Nash}
\label{sec:main}

\begin{theorem}[Quantitative Well-Posedness]
\label{thm:wellposedness}
Under Assumptions~\ref{ass:lipschitz-drift}--\ref{ass:initial}, for any
$\lambda > \tfrac{C_1}{2} + \tfrac{C_2}{2}$, the map $\Phi$ is a strict contraction on
the complete metric space $(\mathcal{M}, \|\cdot\|_{\lambda,T})$:
\begin{equation}
    \|\Phi(\mu^1) - \Phi(\mu^2)\|_{\lambda,T} \le \rho\, \|\mu^1-\mu^2\|_{\lambda,T},
    \qquad \rho = \sqrt{\frac{C_2}{2\lambda - C_1}} < 1.
    \label{eq:contraction}
\end{equation}
Consequently, there exists a unique fixed point $\mu^* = \Phi(\mu^*) \in \mathcal{M}$,
and the corresponding processes $(X^*, p^*, q^*, r^*, \alpha^*)$ solving
\eqref{eq:agent-fbsde} with $\mu=\mu^*$ constitute the unique solution to the
common-noise MFG equilibrium system. Moreover, if all $N$ players in the associated
finite-player game adopt the strategy $\alpha^*$, the resulting strategy profile is an
$\varepsilon_N$-Nash equilibrium with $\varepsilon_N = O(N^{-1/2})$.
\end{theorem}

\begin{proof}[Proof sketch]
The contraction bound \eqref{eq:contraction} follows by taking the weighted supremum
of \eqref{eq:stability-estimate} and bounding the resulting convolution integral,
exactly as in the derivation of \eqref{eq:constants}; see \Cref{app:contraction-proof}
for the full computation. Existence and uniqueness of the fixed point then follow from
the Banach fixed-point theorem, since $(\mathcal{M}, \|\cdot\|_{\lambda,T})$ is
complete (inherited from completeness of $\mathcal{P}_2(\mathbb{R}^d)$ and standard
properties of Skorokhod space; \citealp{billingsley1999}). The $\varepsilon$-Nash
property follows from conditional propagation of chaos
\citep[Vol.~II, Thm.~6.16]{carmonadelarue2018}, applied to bound the distance between
the finite-player empirical measure and $\mu^*$, combined with the Lipschitz stability
of the FBSDE system established above.
\end{proof}

\section{Discussion of the Contraction Constant}
\label{sec:discussion}

The constant $C_1 = 2L_b^2 + 2\|\sigma_0\|_F^2 + C_{\mathrm{reg}}$ in
\eqref{eq:constants} admits a clean decomposition:
\begin{itemize}
    \item $2L_b^2$: the standard Lipschitz contribution from the drift, present even
    without common noise;
    \item $2\|\sigma_0\|_F^2$: the effect of the common noise, entering not through
    the (cancelling) first-order stochastic term but through the dependence of the
    BSDE constant $C_{\mathrm{BSDE}}$ on $\|\sigma_0\|_F$ via the martingale
    representation for $r_t$;
    \item $C_{\mathrm{reg}}$: regularity constants from the Lions derivatives of
    $b, f, g$ and the adjoint BSDE bounds. This term is fully explicit for specific
    models (e.g., the linear-quadratic case of \Cref{sec:example}) and otherwise
    bounded in terms of Lipschitz constants of the Lions derivatives.
\end{itemize}

\paragraph{Limitations.} The estimate relies on absolute continuity and bounded
densities for the stability of optimal transport maps (\Cref{ass:regularity}). In
particle-method implementations, where the empirical measure is discrete, this
assumption requires regularisation (e.g., entropic smoothing); the error introduced
by such regularisation is not quantified here. The contraction constant $\rho$ also
grows with the horizon $T$ through $C_1, C_2$; for long horizons this may require
choosing a large $\lambda$, and under weaker structural conditions the contraction
established here may only hold locally in time.

\section{Illustration: The Linear-Quadratic Benchmark}
\label{sec:example}

\begin{example}[Canonical LQ-MFG]
\label{ex:lq-mfg}
\begin{align*}
    \text{State:}\quad & dX_t = [\kappa(\bar\mu_t - X_t) + \alpha_t]\, dt + \sigma\,
        dW_t + \sigma_0\, dW_t^0, \\
    \text{Running cost:}\quad & f(t,x,\mu,\alpha) = \tfrac{c_1}{2}x^2 +
        \tfrac{\lambda}{2}\alpha^2, \\
    \text{Terminal cost:}\quad & g(x,\mu) = \tfrac{c_2}{2}(x-\bar\mu)^2,
\end{align*}
with $\bar\mu_t = \int x\,\mu_t(dx)$, $\kappa, \lambda, c_1, c_2 > 0$.
\end{example}

For this model, every constant in \Cref{thm:stability,thm:wellposedness} is explicit
in terms of $\kappa, \lambda, c_1, c_2, \sigma, \sigma_0$. The value function is
$V(t,x,\mu) = -\tfrac12 A_t x^2 - B_t x\bar\mu - \tfrac12 C_t\bar\mu^2 - D_t x -
E_t\bar\mu - F_t$, with $(A_t,\ldots,F_t)$ solving a coupled Riccati system derived
from the master equation for this model. With baseline parameters $\kappa=0.5$,
$\lambda=1$, $c_1=2$, $c_2=5$, $\sigma=0.2$, $\sigma_0=0.3$, $T=1$, the Riccati system
is well-posed on $[0,T]$ and yields an explicit contraction constant, verified
numerically by two independent ODE solvers (LSODA and Radau). The full derivation of
this Riccati system, including a re-derivation identifying and resolving two
discrepancies against an earlier working version (a sign difference in the quadratic
coefficient and missing mean-field coupling terms), is provided in a companion
technical note.

\section{Conclusion}

We have established a quantitative, parameter-dependent contraction estimate for the
fixed-point map characterising equilibria of common-noise mean-field games, extending
the qualitative well-posedness theory of \citet{carmonadelarue2018} and
\citet{ahujarenyang2019} with an explicit rate suitable for algorithmic use. Natural
directions for future work include extending the explicit-constant analysis beyond
the linear-quadratic case, quantifying the regularisation error for discrete
(particle-based) approximations, and using the contraction rate derived here as a
design criterion for convergent numerical schemes -- a direction pursued in a
companion paper on numerical approximation of the coupled FBSDE--Fokker--Planck
system.

\bibliographystyle{plainnat}
\begin{thebibliography}{99}

\bibitem[Achdou and Capuzzo-Dolcetta(2010)]{achdoucapuzzo2010}
Achdou, Y. and Capuzzo-Dolcetta, I. (2010).
\newblock Mean field games: numerical methods.
\newblock \emph{SIAM J. Numer. Anal.}, 48(3):1136--1162.

\bibitem[Ahuja et al.(2019)]{ahujarenyang2019}
Ahuja, S., Ren, W., and Yang, T.-W. (2019).
\newblock Forward-backward stochastic differential equations with monotone
functionals and mean field games with common noise.
\newblock \emph{Stochastic Process. Appl.}, 129(10):3859--3892.

\bibitem[Billingsley(1999)]{billingsley1999}
Billingsley, P. (1999).
\newblock \emph{Convergence of Probability Measures}, 2nd edition.
\newblock Wiley.

\bibitem[Cardaliaguet(2010)]{cardaliaguet2010notes}
Cardaliaguet, P. (2010).
\newblock Notes on mean field games.
\newblock Technical report, Universit\'e Paris-Dauphine.

\bibitem[Cardaliaguet et al.(2019)]{cardaliaguetetal2019}
Cardaliaguet, P., Delarue, F., Lasry, J.-M., and Lions, P.-L. (2019).
\newblock \emph{The Master Equation and the Convergence Problem in Mean Field
Games}.
\newblock Princeton University Press.

\bibitem[Cardaliaguet et al.(2025)]{cardaliaguetseegersouganidis2025}
Cardaliaguet, P., Seeger, B., and Souganidis, P. E. (2025).
\newblock Taming under isoperimetry.
\newblock \emph{Stochastic Process. Appl.} arXiv:2311.09003.

\bibitem[Carmona and Delarue(2018)]{carmonadelarue2018}
Carmona, R. and Delarue, F. (2018).
\newblock \emph{Probabilistic Theory of Mean Field Games with Applications I--II}.
\newblock Springer.

\bibitem[Carmona et al.(2013)]{carmona2013probabilistic}
Carmona, R., Delarue, F., and Lachapelle, A. (2013).
\newblock Control of McKean--Vlasov dynamics versus mean field games.
\newblock \emph{Math. Financ. Econ.}, 7(2):131--166.

\bibitem[Carmona and Lauri\`ere(2022)]{carmonalauriere2022}
Carmona, R. and Lauri\`ere, M. (2022).
\newblock Deep learning for mean field games and mean field control with
applications to finance.
\newblock In \emph{Machine Learning and Data Sciences for Financial Markets}.
Cambridge University Press.

\bibitem[El Karoui et al.(1997)]{elkarouietal1997}
El Karoui, N., Peng, S., and Quenez, M. C. (1997).
\newblock Backward stochastic differential equations in finance.
\newblock \emph{Math. Finance}, 7(1):1--71.

\bibitem[Fournier and Guillin(2015)]{fournierguillin2015}
Fournier, N. and Guillin, A. (2015).
\newblock On the rate of convergence in Wasserstein distance of the empirical
measure.
\newblock \emph{Probab. Theory Related Fields}, 162(3--4):707--738.

\bibitem[Huang et al.(2006)]{huangetal2006}
Huang, M., Caines, P. E., and Malham\'e, R. P. (2006).
\newblock Large population stochastic dynamic games: closed-loop McKean--Vlasov
systems and the Nash certainty equivalence principle.
\newblock \emph{Commun. Inf. Syst.}, 6(3):221--252.

\bibitem[Lasry and Lions(2007)]{lasrylions2007}
Lasry, J.-M. and Lions, P.-L. (2007).
\newblock Mean field games.
\newblock \emph{Jpn. J. Math.}, 2(1):229--260.

\bibitem[McKean(1966)]{mckean1966}
McKean, H. P. (1966).
\newblock A class of Markov processes associated with nonlinear parabolic
equations.
\newblock \emph{Proc. Natl. Acad. Sci. USA}, 56(6):1907--1911.

\bibitem[Sznitman(1991)]{sznitman1991}
Sznitman, A.-S. (1991).
\newblock Topics in propagation of chaos.
\newblock In \emph{\'Ecole d'\'Et\'e de Probabilit\'es de Saint-Flour XIX--1989},
pages 165--251. Springer.

\bibitem[Friz et al.(2026)]{frizetal2026}
Friz, P. K., Gasteratos, I., Horst, U., and Theodorakopoulos, S. (2026).
\newblock Mean-field games with rough common noise I: the linear-quadratic case.
\newblock arXiv:2602.19210.

\bibitem[Jackson and Mészáros(2025)]{jacksonmeszaros2025}
Jackson, J. and Mészáros, A. R. (2025).
\newblock Quantitative convergence for displacement monotone mean field games of
control.
\newblock arXiv:2601.07028.

\bibitem[Jackson and Tangpi(2024)]{jacksontangpi2024}
Jackson, J. and Tangpi, L. (2024).
\newblock Quantitative convergence for displacement monotone mean field games with
controlled volatility.
\newblock \emph{Math. Oper. Res.}, 49(4):2527--2564.

\bibitem[Peng and Wu(1999)]{peng1999}
Peng, S. and Wu, Z. (1999).
\newblock Fully coupled forward-backward stochastic differential equations and
applications to optimal control.
\newblock \emph{SIAM J. Control Optim.}, 37(3):825--843.

\bibitem[Zhang(2017)]{zhang2017}
Zhang, J. (2017).
\newblock \emph{Backward Stochastic Differential Equations}.
\newblock Springer.

\end{thebibliography}

\appendix

\section{Algebraic Details: Closing the Gronwall Inequality}
\label{app:algebra}

This appendix gives the derivation of \eqref{eq:gronwall-differential} from
\eqref{eq:gronwall-setup} and \eqref{eq:bsde-estimate}, together with the control
bound \eqref{eq:control-lipschitz}.

\paragraph{Step 1: bound the control difference in the $X$-inequality.} By
\eqref{eq:control-lipschitz} and the elementary inequality
$(a+b+c)^2 \le 3(a^2+b^2+c^2)$,
\[
    |\alpha_t^1-\alpha_t^2|^2 \le 3C_\alpha^2\big(|\Delta X_t|^2 + |\Delta p_t|^2 +
    W_2(\mu_t^1,\mu_t^2)^2\big).
\]
Substituting into \eqref{eq:gronwall-setup},
\begin{equation}
    d|\Delta X_t|^2 \le \Big[\big(2L_b+L_b^2+3L_bC_\alpha^2\big)|\Delta X_t|^2 +
    3L_bC_\alpha^2\,|\Delta p_t|^2 + \big(L_b+3L_bC_\alpha^2\big)
    W_2(\mu_t^1,\mu_t^2)^2\Big]\, dt.
    \label{eq:app-x-bound}
\end{equation}

\paragraph{Step 2: bound the control difference in the BSDE estimate.} Applying
\eqref{eq:bsde-estimate} on the sub-interval $[0,t]$ (valid since $t \le T$ is
arbitrary and the same BSDE comparison argument applies verbatim on any
sub-interval), and substituting the same bound on $|\alpha^1_s-\alpha^2_s|^2$:
\begin{align}
    \mathbb{E}\Big[\sup_{s\le t}|\Delta p_s|^2\Big] &\le C_{\mathrm{BSDE}}\,
    \mathbb{E}\big[|\Delta X_t|^2\big] + C_{\mathrm{BSDE}}(1+3C_\alpha^2)
    \int_0^t \mathbb{E}\big[|\Delta X_s|^2\big]\, ds \notag \\
    &\qquad + 3C_{\mathrm{BSDE}} C_\alpha^2 \int_0^t
    \mathbb{E}\big[|\Delta p_s|^2\big]\, ds + C_{\mathrm{BSDE}}(1+3C_\alpha^2)
    \int_0^t \mathbb{E}\big[W_2(\mu_s^1,\mu_s^2)^2\big]\, ds.
    \label{eq:app-p-bound}
\end{align}

\paragraph{Step 3: a smallness condition closes the estimate.} The term
$3C_{\mathrm{BSDE}}C_\alpha^2 \int_0^t \mathbb{E}[|\Delta p_s|^2]\,ds$ on the
right-hand side of \eqref{eq:app-p-bound} involves the same quantity bounded on
the left. We assume, as part of the standing regularity hypotheses
(\Cref{ass:regularity}), that the model parameters satisfy
\begin{equation}
    3\,C_{\mathrm{BSDE}}\,C_\alpha^2 < 1;
    \label{eq:smallness-condition}
\end{equation}
this is a mild restriction on the size of $C_\alpha$ (control sensitivity) relative
to $C_{\mathrm{BSDE}}$ (adjoint-process sensitivity), both of which are themselves
determined by the Lipschitz and concavity constants of the model, and can always be
arranged by working on a sufficiently short time interval or under a sufficiently
strong concavity constant $\lambda_H$ in \Cref{ass:concavity}. Under
\eqref{eq:smallness-condition}, a standard integral form of Gronwall's lemma applied
to \eqref{eq:app-p-bound} (treating $\mathbb{E}[\sup_{s\le t}|\Delta p_s|^2]$ as the
unknown and absorbing the self-referential term via
\eqref{eq:smallness-condition}) yields
\begin{equation}
    \mathbb{E}\big[|\Delta p_t|^2\big] \le \frac{C_{\mathrm{BSDE}}}{1-3C_{\mathrm{BSDE}}C_\alpha^2}
    \left(\mathbb{E}\big[|\Delta X_t|^2\big] + (1+3C_\alpha^2)\int_0^t
    \mathbb{E}\big[|\Delta X_s|^2 + W_2(\mu_s^1,\mu_s^2)^2\big]\, ds\right).
    \label{eq:app-p-closed}
\end{equation}

\paragraph{Step 4: combine.} Adding \eqref{eq:app-x-bound} (in integrated form) and
\eqref{eq:app-p-closed}, and setting
$y_t := \mathbb{E}[|\Delta X_t|^2] + \mathbb{E}[|\Delta p_t|^2]$, every term on the
right is a constant multiple of either $\mathbb{E}[|\Delta X_s|^2]$,
$\mathbb{E}[|\Delta p_s|^2]$ (hence of $y_s$), or $\mathbb{E}[W_2(\mu_s^1,\mu_s^2)^2]$,
for $s \le t$. Collecting terms and denoting
$\widetilde{C}_{\mathrm{BSDE}} := C_{\mathrm{BSDE}}/(1-3C_{\mathrm{BSDE}}C_\alpha^2)$,
one obtains, after routine algebra,
\[
    \frac{d}{dt}\, y_t \;\le\; C_1\, y_t \;+\; C_2\, \mathbb{E}\big[W_2(\mu_t^1,\mu_t^2)^2\big],
\]
where
\[
    C_1 = 2L_b + L_b^2 + 3L_bC_\alpha^2 + \widetilde{C}_{\mathrm{BSDE}}(1+3C_\alpha^2),
    \qquad
    C_2 = L_b + 3L_bC_\alpha^2 + \widetilde{C}_{\mathrm{BSDE}}(1+3C_\alpha^2).
\]
This is \eqref{eq:gronwall-differential}. The compact expression
$C_1 = 2L_b^2+2\|\sigma_0\|_F^2+C_{\mathrm{reg}}$ quoted in \eqref{eq:constants}
absorbs the constant $\widetilde{C}_{\mathrm{BSDE}}(1+3C_\alpha^2)$ and the
$3L_bC_\alpha^2$ term into $C_{\mathrm{reg}}$, and isolates the explicit
$\|\sigma_0\|_F^2$-dependence that enters through $C_{\mathrm{BSDE}}$ (and hence
through $\widetilde{C}_{\mathrm{BSDE}}$) via the martingale representation for
$r_t$ in the adjoint BSDE, as discussed in \Cref{sec:discussion}.

\begin{remark}[On the smallness condition \eqref{eq:smallness-condition}]
This condition does not appear explicitly in the informal derivation sketched in
the main text; making it explicit here is, in the authors' view, necessary for a
fully rigorous closure of the Gronwall argument, and we flag it as such rather than
leaving it implicit. It is a standing structural assumption on the model (via
$C_\alpha, C_{\mathrm{BSDE}}$), not an additional constraint imposed after the fact;
in the linear-quadratic case of \Cref{sec:example} it can be checked directly from
the explicit Riccati coefficients.
\end{remark}

\section{Proof of the Contraction Bound}
\label{app:contraction-proof}

This appendix derives \eqref{eq:contraction} from \eqref{eq:stability-estimate}.
Multiplying both sides of \eqref{eq:stability-estimate} by $e^{-2\lambda t}$ and
taking the supremum over $t \in [0,T]$:
\begin{align*}
    \|\Phi(\mu^1)-\Phi(\mu^2)\|_{\lambda,T}^2 &= \sup_{t\le T} e^{-2\lambda t}\,
        \mathbb{E}\big[W_2(\Phi(\mu^1)_t,\Phi(\mu^2)_t)^2\big] \\
    &\le C_2 \sup_{t\le T} e^{-2\lambda t} \int_0^t e^{C_1(t-s)}\,
        \mathbb{E}\big[W_2(\mu_s^1,\mu_s^2)^2\big]\, ds.
\end{align*}
Writing the integrand as
$e^{-2\lambda(t-s)}\, e^{(C_1-2\lambda)s} \cdot e^{-2\lambda s}\,
\mathbb{E}[W_2(\mu_s^1,\mu_s^2)^2]$
and using $e^{-2\lambda s}\,\mathbb{E}[W_2(\mu_s^1,\mu_s^2)^2] \le
\|\mu^1-\mu^2\|_{\lambda,T}^2$ for every $s\le T$, together with the change of
variable $u = t-s$:
\[
    \|\Phi(\mu^1)-\Phi(\mu^2)\|_{\lambda,T}^2 \le C_2\left(\sup_{t\le T}
    \int_0^t e^{-(2\lambda-C_1)u}\, du\right) \|\mu^1-\mu^2\|_{\lambda,T}^2.
\]
For $\lambda > C_1/2$, the integral is increasing in $t$ and bounded by its value
at $t=T$:
\[
    I(\lambda) := C_2\int_0^T e^{-(2\lambda-C_1)u}\, du =
    \frac{C_2}{2\lambda-C_1}\big(1-e^{-(2\lambda-C_1)T}\big) \le
    \frac{C_2}{2\lambda-C_1}.
\]
Choosing $\lambda > \tfrac{C_1}{2}+\tfrac{C_2}{2}$ ensures $2\lambda - C_1 > C_2$,
hence $I(\lambda) < 1$, and $\rho := \sqrt{I(\lambda)} = \sqrt{C_2/(2\lambda-C_1)}
< 1$. This gives
\[
    \|\Phi(\mu^1)-\Phi(\mu^2)\|_{\lambda,T} \le \rho\, \|\mu^1-\mu^2\|_{\lambda,T},
\]
which is \eqref{eq:contraction}. Since $(\mathcal{M}, \|\cdot\|_{\lambda,T})$ is a
complete metric space (\Cref{sec:setup}), the Banach fixed-point theorem yields a
unique $\mu^* \in \mathcal{M}$ with $\Phi(\mu^*)=\mu^*$, completing the proof of
existence and uniqueness in \Cref{thm:wellposedness}. \qed

\end{document}
